\documentclass[12pt]{article}
\usepackage[T1]{fontenc}
\usepackage[utf8]{inputenc}
\usepackage{lmodern}
\usepackage{textcomp}
\usepackage{enumitem}
\usepackage{geometry}
\geometry{margin=1in}
\begin{document}
\begin{center}
Answer Key - Constitutional Law - Graduate Level Assessment 6 \
\small (Answers are ordered exactly as in the question paper.)
\end{center}

\section*{Multiple Choice}
\begin{enumerate}[label=\arabic*.]
\item \textbf{Answer:} B
\\ \textit{Options:} A) Indisputable facts.; B) Facts that have been asserted by individual political organizations.; C) Facts recognized to be true by common knowledge.; D) Facts capable of scientific verification.
\\ \textit{Note:} Facts asserted by individual political organizations are not appropriate for judicial notice because they lack the requisite reliability and universal acceptance that characterize facts suitable for judicial notice in a court of law.
\item \textbf{Answer:} D
\\ \textit{Options:} A) F.A.S. Contract; B) Destination contract; C) Shipment contract; D) F.O.B. Contract
\\ \textit{Note:} An F.O.B. (Free On Board) contract specifies the delivery point as a critical element of the agreement, indicating where the ownership and responsibility for the goods transfer from the seller to the buyer.
\item \textbf{Answer:} B
\\ \textit{Options:} A) Be made by a merchant to a merchant.; B) Be contained in a signed writing which gives assurance that the offer will be held open.; C) State the period of time for which it is irrevocable.; D) Not be contained in a form supplied by the offeror.
\\ \textit{Note:} The correct answer is based on the Uniform Commercial Code's requirement that an irrevocable offer must be made in a signed writing that provides assurance of its open status, thereby ensuring clarity and commitment in commercial transactions.
\item \textbf{Answer:} D
\\ \textit{Options:} A) Pay just compensation for the property; B) Terminate the regulation and pay damages that occurred while regulation was in effect; C) Pay double damages if not addressed in 30 days.; D) a or b
\\ \textit{Note:} The correct answer is "a or b" because, under the Fifth Amendment's Takings Clause, if a regulation constitutes a taking of private property for public use, the government is required to either compensate the property owner or modify the regulation to avoid the taking.
\item \textbf{Answer:} B
\\ \textit{Options:} A) Executive agreements, U.S. Constitution, treaties and federal statutes, state law; B) U.S constitution, treaties and federal statutes, executive agreements, state law; C) Treaties and federal statutes, US Constitution, executive agreements, state law; D) U.S constitution, executive agreements, treaties and federal statutes, state law
\\ \textit{Note:} The correct answer reflects the established legal hierarchy in the United States, where the U.S. Constitution is the supreme law of the land, followed by treaties and federal statutes, then executive agreements, and finally state law, ensuring a structured framework for legal authority and governance.
\end{enumerate}

\section*{True / False}
\begin{enumerate}[label=\arabic*.]
\setcounter{enumi}{5}
\item \textbf{Answer:} False
\\ \textit{Options:} A) True; B) False
\\ \textit{Note:} Nominal damages are meant to recognize a legal wrong without providing substantial financial compensation, as they typically award a small sum to confirm that a breach occurred rather than to cover lost profits.
\item \textbf{Answer:} False
\\ \textit{Options:} A) True; B) False
\\ \textit{Note:} The name 'FBI Consultants, Incorporated' is misleading and could imply a false affiliation or endorsement by the Federal Bureau of Investigation, which violates trademark and false advertising laws under constitutional principles.
\item \textbf{Answer:} False
\\ \textit{Options:} A) True; B) False
\\ \textit{Note:} The Rule in Shelley's Case invalidates a remainder interest in the heirs of the grantor, causing it to merge with the preceding life estate and resulting in a fee simple estate for the life tenant, rather than allowing a valid remainder that would prevent reversion to the grantor.
\item \textbf{Answer:} False
\\ \textit{Options:} A) True; B) False
\\ \textit{Note:} The statement is false because a contract can be discharged due to frustration of purpose even if the frustrating event was foreseeable at the time of contract formation, as long as the event fundamentally changes the nature of the contractual obligations.
\item \textbf{Answer:} False
\\ \textit{Options:} A) True; B) False
\\ \textit{Note:} The statement is false because even when a law appears discriminatory on its face, courts often require evidence of discriminatory intent to establish a violation of equal protection principles under constitutional law.
\end{enumerate}

\section*{Long Form Response}
\begin{enumerate}[label=\arabic*.]
\setcounter{enumi}{10}
\item \textbf{Answer:} The President's power to appoint various officials with the advice and consent of the Senate plays a crucial role in maintaining the balance of power within the federal government. This process ensures that while the President can select key figures in the executive branch, such as cabinet members and federal judges, these appointments require Senate approval, thus providing a check on executive authority. This collaborative aspect fosters accountability and encourages the consideration of diverse perspectives, which can mitigate the risk of unilateral decision-making. Additionally, the Senate's role in confirming appointments reinforces the principle of separation of powers by ensuring that no single branch can dominate the federal government. Ultimately, this appointment process exemplifies the interplay between different branches of government, highlighting the importance of cooperation and oversight in a functioning democracy.
\\ \textit{Note:} correct because it highlights the essential role of the Senate's advice and consent in checking presidential power and promoting accountability within the executive branch, thereby preserving the balance of power in the federal government.
\item \textbf{Answer:} Slander, a form of defamation, requires that a statement be false, spoken, and damaging to the reputation of an individual. For a statement to be considered slanderous, it typically must meet the criteria of being made with actual malice or negligence regarding its truthfulness, particularly when the subject is a public figure. However, in cases of slander per se, certain statements are deemed inherently damaging, such as those accusing someone of a crime or suggesting unchastity, which do not require proof of harm. This raises the important implication that publication is not necessary for slander per se, as the gravity of the statement itself is sufficient to warrant legal action. Thus, the focus shifts from the act of publication to the nature of the statement, allowing for immediate legal recourse in cases of egregious claims.
\\ \textit{Note:} correct because it accurately outlines the fundamental requirements for a statement to qualify as slander, including the necessity for the statement to be false, spoken, and damaging, while also addressing the concept of slander per se, where certain statements are considered damaging without the need for proof of publication.
\end{enumerate}

\vfill
\noindent\textit{End of Answer Key}
\end{document}